\section{Related work}\label{sec:related}

\subsection{Open-vocabulary detection}

Open-vocabulary detection replaces a fixed classifier with a query: a category is described
at inference time and the detector scores regions against that description. The two families
that matter here differ in what the query is made of. Region--text alignment methods score a
region against a word embedding learned from image--caption data, either by contrasting
whole images with sentences and transferring the classifier to regions
\citep{radford2021learning} or by supervising the region--phrase alignment directly
\citep{zhong2022regionclip,yao2022detclip,minderer2023scaling}. Grounded detectors instead
fuse text tokens with image features inside the detection head, so that the classification
logit of an object query is computed against the token embeddings of the phrase. The
transformer decoder that made this possible was designed for closed-set detection
\citep{carion2020end}; MDETR conditioned it on a whole sentence \citep{kamath2021mdetr}, GLIP
recast detection as phrase grounding and scaled the grounding supervision
\citep{li2022grounded}, Grounding DINO added a cross-modality fusion stage to a DINO-style
detector \citep{liu2024grounding}, and YOLO-World and OV-DINO reorganised the same interface
for real-time inference and for a unified language-aware fusion
\citep{cheng2024yolo,wang2024ovdino}. A parallel line lets the query contain boxes as well as
words, in the form of exemplars drawn on the image
\citep{xu2023multi,jiang2023trex,jiang2024trex2}, and segmentation at any granularity shares
the interface \citep{li2024segment}. In every one of these designs the score that the
user reads is a similarity between a region and a representation of the \emph{query}, not a
posterior over a fixed label set. That is the property the rest of this paper turns on.

The literature on the query has concentrated on accuracy: which prompt template gives the
best average precision, whether prompt tuning helps \citep{zhou2022learning}, how to ensemble
prompt embeddings \citep{radford2021learning}, how to fold the vocabulary into the weights so
that a text encoder is not needed at inference time \citep{cheng2024yolo}. The question we
ask is different, and orthogonal to all of them: given that the score is a function of the
wording, what happens to the \emph{meaning} of the score when the wording changes, and what
does a threshold fitted under one wording mean under another?

\subsection{What a detection score is a probability of}

That a detector's score is not a probability is not a new observation.
\citet{hall2020probabilistic} define probabilistic object detection as the task of reporting a
detection together with its spatial and semantic uncertainty, propose an evaluation measure
without arbitrary score thresholds, and find that conventional detectors are spatially
overconfident. Monte Carlo dropout has been used in the same spirit to extract label
uncertainty from a detector and to reject objects from outside the training set
\citep{miller2018dropout}. Both treat the label as given: the uncertainty is over \emph{where}
the object is and \emph{which} of the known classes it belongs to, and the quantity that is
miscalibrated is a property of the image and the network. Once the label is a sentence typed
at inference time, the object of uncertainty changes. The same region scored against
\emph{dog} and against \emph{a photo of a dog} is two different queries, and the two scores are
two different claims about the same pixels. This paper is about the second kind of
variability, which the probabilistic detection literature does not have to consider because
its label set is fixed.

\subsection{Calibration of detectors}

Post-hoc calibration maps a score to a probability using held-out data. The classical options
are parametric, Platt scaling \citep{platt1999probabilistic} and its per-group generalisation,
and non-parametric, histogram binning and isotonic regression \citep{zadrozny2002transforming};
temperature scaling \citep{guo2017calibration} is the one-parameter version that leaves
accuracy untouched because it is monotone. Calibrated regression
\citep{kuleshov2018accurate} generalises the idea to continuous outputs, and a body of work
examines when and how these estimators fail
\citep{vaicenavicius2019evaluating,nixon2019measuring,wang2023calibration}.

Two assumptions of that literature are violated in our setting. First, calibration is usually
evaluated as an average over the whole score range, which for a detector is dominated by the
boxes nobody keeps; we therefore report the error inside the band where review actually
happens. Second, and more importantly for this paper, the calibrated object is assumed to be
the model's output on a fixed input distribution. When the input includes a user-typed phrase,
the distribution that matters is indexed by the phrase, and a map fitted on the marginal over
phrasings is not the right map for any of them.

The closest precedent lives in a different modality, and it is close enough that the
difference is worth stating precisely. \citet{zhao2021calibrate} observed that the few-shot
accuracy of a language model swings with the prompt format and traced the swing to a bias
towards particular answers; they estimate that bias with a content-free input such as
``N/A'' and fit, per prompt, an affine map that makes the content-free prediction uniform.
The premise --- that the object to calibrate is the pair (model, prompt) rather than the model
alone --- is the one we adopt, and the affine form is the same. What differs is how the map is
obtained and what it is asked to deliver. A detector has no content-free input: there is no
image whose detection scores would reveal the bias of the phrase \emph{a photo of a}, so the
map cannot be measured at deployment from unlabelled probes and has to be predicted from the
phrase itself. And the target here is not an aggregate accuracy but the calibration of
individual detections inside an operating band, together with the ordering that a threshold
acts on. Section~\ref{sec:exp-fdr} shows that the second of these is out of reach for
any recalibration, contextual or otherwise, which is why the phrase-conditioned map is only
half of our proposal and the threshold protocol is the other half.

\subsection{Conformal selection and risk control}

Turning a score into a decision with a guarantee is a separate problem from making the score a
probability. Selective prediction trades coverage for accuracy
\citep{geifman2017selective,geifman2019selectivenet}, and conformal prediction supplies
distribution-free coverage guarantees \citep{angelopoulos2023conformal}. Risk-controlling
prediction sets \citep{bates2021distribution} generalise the guarantee from coverage to an
arbitrary monotone loss calibrated on a holdout set, and conformalised selection
\citep{jin2023selection} is the false-discovery-rate case of that framework, which is the
quantity a user of a detection tool cares about: of the boxes I keep, what fraction is wrong.
The procedure needs a labelled calibration sample and an exchangeability assumption between
calibration and test. In a deployment where the user types an arbitrary phrase,
exchangeability holds comfortably within a phrasing and not at all across phrasings. We adopt
the procedure and then test that assumption directly, and the answer is negative: a threshold
fitted on a mixture of wordings does not deliver its target on the wordings the mixture
contains, and no monotone recalibration of the score changes that, because the target depends
on the ordering of the boxes rather than on the scale. Requiring the constraint to hold
\emph{inside} each wording does deliver it, at a cost in kept boxes that is smaller when the
scores of the wordings already agree --- which is the sense in which conditioning the phrase
pays for itself.

\subsection{Sensitivity to the prompt}

That a prompt's wording moves the output is by now well documented, mostly for language
models. The \emph{format} of a prompt alone --- separators, casing, spacing --- produces
accuracy spreads comparable to the differences between models, and the spread is stable
enough to be measured by a sensitivity metric computed over a set of format variations
\citep{sclar2024quantifying}. The \emph{order} of few-shot examples does the same: some
permutations are near state of the art and others near chance, and a good permutation for one
model does not transfer to another \citep{lu2022fantastically}. A stronger result is that many
intentionally irrelevant or even misleading prompts work as well as carefully written ones,
which undermines the reading that a prompt is understood as an instruction
\citep{webson2022prompts}. On the vision--language side the template matters enough that a
prompt is learned rather than written \citep{zhou2022learning}, and CLIP already ensembles
several hand-written templates because no single one is best \citep{radford2021learning}.

All of this work measures the effect on accuracy: a dataset-level scalar, in a classification
task where the label set is fixed and every input has one correct answer. The effect we study
is of the same origin --- the wording is arbitrary from the model's point of view, which is
the lesson of \citet{webson2022prompts} --- but it is measured somewhere else. We measure it on
the probability assigned to an \emph{individual detection}, and the quantity at stake is not
accuracy but the meaning of a threshold: a detector whose average precision barely moves can
still have an operating point that means something entirely different under a different
sentence, and that is the case a deployment notices.

\subsection{What is new here}

Calibration, conformal selection and prompt sensitivity are mature as separate literatures,
and the idea of calibrating per prompt exists in another modality. To our knowledge this is
the first study that (i) measures the phrasing-induced shift as a property of the individual
detection rather than of the dataset-level average precision, (ii) predicts a calibrator's
parameters from the phrase representation, so that a wording not seen during fitting is
calibrated without a deployment-time probe, (iii) separates what calibration can buy from
what it cannot --- it moves the scale, so it repairs the band, and it cannot move the ordering,
so it cannot rescue a shared threshold --- and locates the coverage gain of a
phrase-conditioned score in the threshold-selection protocol rather than in the calibrator
alone, and (iv) evaluates the five wordings as five separate runs instead of pooling them,
which is what a deployment actually sees.
